\section{One potential and a common temperature singularity}\label{sec:construction}
We now replace each ideal quadratic attraction by an actual radial kernel and
retain a common ball of supporting slopes. Throughout this section, $t$ is
fixed at a value supplied by Proposition~\ref{prop:initial-triple}; it will
never tend to infinity again. All pressure estimates therefore use the
fixed-$t$ constants proved above. The ideal pressures carry the ball of
supporting slopes, while their physical approximations become the references
for the next step.

\subsection{Choose the next mass before the current range}
Recall the objects supplied by Proposition~\ref{prop:initial-triple}:
$A=3/(2p)$, $q_0=Q_{\psi_t}$, $P_1=T_Aq_0$, the point
$\lambda_*=(1,-3/4)$, a closed ball $W$ of radius $r_0=1/4$ about
$\lambda_*$, three small neighborhoods $U_i$ of the slope vectors $v_i$, and
$g=(v_1+v_2+v_3)/3$. Every triangle with one vertex in each $U_i$ contains
the same ball $\overline B(g,\gamma)$, for some $\gamma>0$.
The initial groups satisfy
$K_i^1\subset\overline B(v_i,\tau/2)$, whereas
$U_i=\overline B(v_i,4\tau)$.
Put $\sigma_1=A$. Choose positive sequences $\zeta_n$ and $\ell_n$ such that
\begin{equation}\label{eq:budgets}
 \sum_{n\ge1}\zeta_n<\tau,
 \qquad \sum_{n\ge1}\ell_n<r_0/16.
\end{equation}
The quantities $\ell_n$ bound displacements in the $(\beta,h)$ plane; they are
unrelated to the side lengths of particle boxes.

Suppose that the finite-range potential and the next ideal pressure have
already been defined by
\begin{align*}
 \varphi_{n-1}&=\psi_t-t\sum_{k<n}\sigma_k w_{R_k},
 &q_{n-1}&=Q_{\varphi_{n-1}},
 &P_n&=T_{\sigma_n}q_{n-1}.
\end{align*}
The induction maintains a point $\lambda_n\in\operatorname{int}W$ and
nonempty compact convex sets $K_i^n$ such that
\begin{equation}\label{eq:inductive-slopes}
 g\in\partial P_n(\lambda_n)
   =\conv\bigcup_{i=1}^3K_i^n,
 \qquad
 K_i^n\subset\overline B\left(v_i,\frac\tau2+
                          \sum_{k<n}\zeta_k\right),
\end{equation}
and every point in $K_i^n$ has nonnegative second coordinate. The geometric
choice of the $U_i$ gives
\begin{equation}\label{eq:common-ball}
 \overline B(g,\gamma)\subset\partial P_n(\lambda_n).
\end{equation}
The enlarged groups $K_i^n+\overline B(0,\zeta_n)$ lie in $U_i$;
their density projections are disjoint, and the convex hull of any two of
them does not contain $g$.

Apply Lemma~\ref{lem:persistence} to $P_n$, $\lambda_n$, and tolerance
$\zeta_n$. Denote its constants by $S_n$ and $a_n^*>0$. First choose
\begin{equation}\label{eq:choose-mass}
 0<\sigma_{n+1}<
 \min\left\{a_n^*,\frac{\ell_n}{\max\{1,S_n\}},2^{-n-1}\right\}.
\end{equation}
Only after this choice, select $R_n>4$, with $R_n>R_{n-1}+1$ for $n>1$,
so large that Proposition~\ref{prop:kac}, applied to the fixed finite-range
reference $\varphi_{n-1}$, gives
\begin{align}
 \varphi_n&=\varphi_{n-1}-t\sigma_nw_{R_n},
 &q_n&=Q_{\varphi_n},\notag\\
 \|q_n-P_n\|_W&\le\min\{\sigma_{n+1}^2,2^{-n}\}.
 \label{eq:choose-range}
\end{align}
The physical class is respected because
$\sum_{k\ge2}\sigma_k<1$, so every partial sum of the masses is at most
$A+1$. Lemma~\ref{lem:persistence} now applies to the physical pressure
$q_n$ with perturbation size $a=\sigma_{n+1}$. It produces
$P_{n+1}=T_{\sigma_{n+1}}q_n$, a point $\lambda_{n+1}$, and groups
$K_i^{n+1}$ satisfying
\[
 |\lambda_{n+1}-\lambda_n|\le\ell_n,
 \qquad K_i^{n+1}\subset K_i^n+\overline B(0,\zeta_n).
\]
This verifies \eqref{eq:inductive-slopes} at the next stage.
The initial location and \eqref{eq:budgets} keep every $\lambda_n$ in
$B(\lambda_*,r_0/8)$. Thus the construction continues indefinitely.

The order in \eqref{eq:choose-mass}--\eqref{eq:choose-range} is essential.
The next attractive mass sets the required precision for approximating the
current ideal pressure. No uniform convergence rate for the Kac limit as the
reference range changes is needed: that reference is fixed before each new
range is chosen.

\subsection{The final radial interaction}
Define
\begin{equation}\label{eq:final-potential}
 \Phi(x)=\psi_t(x)-t\sum_{j=1}^{\infty}\sigma_jw_{R_j}(x).
\end{equation}
Since $R_j\ge1$ and $\sum_j\sigma_j<\infty$, the attractive series
converges uniformly and in $L^1(\R^3)$. It is radial and nonnegative before
subtraction. Consequently $\Phi$ is Borel measurable, finite and locally
bounded away from the origin, and has the same divergence at the origin as
$\psi_t$. The core vanishes for $|x|\ge1$, so the $L^1$ convergence gives
\[
 4\pi\int_1^\infty r^2|\phi(r)|\dd r
 =\int_{|x|\ge1}|\Phi(x)|\dd x<\infty.
\]
The compact support of $w$ also gives, for $r\ge1$,
\[
 r^3|\phi(r)|\le t\|w\|_\infty\sum_{R_j\ge r}\sigma_j
 \longrightarrow0.
\]
This estimates the tail without imposing any power-law relation between
$\sigma_j$ and $R_j$. No bound of order $r^{-3-\varepsilon}$ for a fixed
$\varepsilon>0$ is established by these choices.
Because $R_1>4$ and $w$ is continuous and strictly positive on the ball of
radius $1/2$, its first attraction alone has a positive lower bound on
$1\le|x|\le2$. As $\sigma_1=A>0$,
\[
 \Phi(x)\le -tA R_1^{-3}
            \min_{|y|\le 2/R_1}w(y)<0
 \qquad(1\le|x|\le2).
\]
Finally, Proposition~\ref{prop:stability} applies to the infinite sum and
gives $U_N^\Phi\ge-B_0N$. Thus \eqref{eq:final-potential} satisfies every
admissibility condition in Section~\ref{sec:introduction}.

Write $Q_\infty=Q_\Phi$ and $H_\infty=H_\Phi$ for the actual pressure and
canonical limit of this potential. Their existence, finiteness, and duality
follow from Proposition~\ref{prop:thermodynamics}. The range-uniform tail
estimate, Proposition~\ref{prop:tail}, and \eqref{eq:choose-range} imply
\begin{equation}\label{eq:physical-pressure-limit}
 \|Q_\infty-P_n\|_W
 \le C_W\sum_{k>n}\sigma_k+\min\{\sigma_{n+1}^2,2^{-n}\}
 \longrightarrow0.
\end{equation}
This estimate identifies the limiting function with the pressure of the
single potential \eqref{eq:final-potential}.

\subsection{Pass the ball of slopes to the physical pressure}
The summable displacement bounds imply
$\lambda_n\to\lambda_c=(\beta_c,h_c)$, with
$\lambda_c\in B(\lambda_*,r_0/8)$.
Fix any $v\in\overline B(g,\gamma)$.
By \eqref{eq:common-ball}, for every $\lambda\in\operatorname{int}W$,
\[
 P_n(\lambda)\ge P_n(\lambda_n)+v\cdot(\lambda-\lambda_n).
\]
Letting $n\to\infty$ in this inequality using
\eqref{eq:physical-pressure-limit} gives a supporting inequality for
$Q_\infty$ at $\lambda_c$ on $\operatorname{int}W$.
Such a local support is global on the convex domain
$\Omega=(0,\infty)\times\R$. Indeed, for any $\lambda\in\Omega$, a short
initial portion of the segment from $\lambda_c$ to $\lambda$ lies in
$\operatorname{int}W$; combining its supporting inequality with convexity
of $Q_\infty$ gives the inequality at $\lambda$. Hence
\begin{equation}\label{eq:limit-ball}
 \overline B(g,\gamma)\subset\partial Q_\infty(\lambda_c).
\end{equation}
The same ball, rather than singularity alone, is the information that survives
all finite-range approximations.

\subsection{A common critical point for an open interval of densities}
The center density is determined by the original triangle:
\begin{equation}\label{eq:fixed-density}
 \rho_*=g_2=\frac{0+p+4p}{3}=\frac{5p}{3}>0.
\end{equation}
Set
\[
 \delta_\rho=\min\{\gamma/2,\rho_*/2\},\qquad
 I=(\rho_*-\delta_\rho,\rho_*+\delta_\rho).
\]
This nonempty interval lies in $(0,\infty)$ and is independent of inverse
temperature. Fix any $\rho\in I$. The vectors
\[
 w_-=(g_1-\gamma/2,\rho),\qquad
 w_+=(g_1+\gamma/2,\rho)
\]
have distance less than $\gamma$ from $g$, so both belong to
$\partial Q_\infty(\beta_c,h_c)$ by \eqref{eq:limit-ball}.
Figure~\ref{fig:fixed-density} shows the central slice; the same argument
applies to every slice at a density in $I$.

\begin{figure}[hbp]
\centering
\begin{tikzpicture}[scale=1.05,>=Stealth]
 \draw[->] (-0.4,0) -- (5.4,0) node[right] {$s$};
 \draw[->] (0,-0.2) -- (0,3.7) node[above] {$d$};
 \fill[blue!8] (2.8,2) circle (1.15);
 \draw[blue!65!black] (2.8,2) circle (1.15);
 \draw[dashed] (0,2) node[left] {$\rho=g_2$} -- (5.0,2);
 \fill (2.8,2) circle (1.4pt) node[below] {$g$};
 \fill (2.22,2) circle (2pt) node[above left] {$w_-$};
 \fill (3.38,2) circle (2pt) node[above right] {$w_+$};
 \node[align=center] at (2.8,0.45) {same density, distinct slopes};
 \node at (2.8,3.5) {$\overline B(g,\gamma)\subset\partial Q_\infty(\beta_c,h_c)$};
\end{tikzpicture}
\caption{The horizontal slice of the preserved subgradient ball gives two
supporting vectors with the same density coordinate. Partial conjugacy turns
their distinct first coordinates into distinct canonical supergradients in
inverse temperature.}\label{fig:fixed-density}
\end{figure}

To obtain the temperature jump, take any $(s,\rho)$ in that subdifferential.
Restricting the
supporting inequality to $\beta=\beta_c$ gives Fenchel equality
\[
 H_\infty(\beta_c,\rho)=h_c\rho-Q_\infty(\beta_c,h_c).
\]
The full inequality, at an arbitrary $(\beta,h)\in\Omega$, rearranges to
\[
 h\rho-Q_\infty(\beta,h)
 \le H_\infty(\beta_c,\rho)-s(\beta-\beta_c).
\]
Taking the supremum over $h$ yields
\begin{equation}\label{eq:canonical-support}
 H_\infty(\beta,\rho)
 \le H_\infty(\beta_c,\rho)-s(\beta-\beta_c),\qquad \beta>0.
\end{equation}
Thus $-s$ is a supergradient of the finite concave function
$H_\infty(\cdot,\rho)$. For a finite concave function at an interior point,
every supergradient lies between the right and left derivatives. Applying
this to $w_-$ and $w_+$ gives
\begin{equation}\label{eq:derivative-jump}
 D_\beta^-H_\infty(\beta_c,\rho)
 -D_\beta^+H_\infty(\beta_c,\rho)\ge\gamma>0.
\end{equation}
Both derivatives are finite because $\beta_c$ is interior to the interval
of finiteness of this concave function.

The normalization \eqref{eq:free-energy-normalization} gives
$f_\Phi(\beta,\rho)=tH_\infty(\beta,\rho)/\beta$. Continuity and the
one-sided product rule therefore give the explicit uniform bound
\[
 D_\beta^- f_\Phi(\beta_c,\rho)
 -D_\beta^+ f_\Phi(\beta_c,\rho)
 =\frac{t}{\beta_c}\left(
 D_\beta^-H_\infty(\beta_c,\rho)
 -D_\beta^+H_\infty(\beta_c,\rho)\right)
 \ge\frac{t\gamma}{\beta_c}>0.
\]
The all-sequence canonical limit supplies
\eqref{eq:free-energy-definition} with $N_L=\lfloor\rho L^3\rfloor$
for every $\beta>0$ and every $\rho>0$. The location of $\lambda_c$
gives $\beta_c\in(1/2,3/2)$. Since the choice of $\rho\in I$ was arbitrary,
this proves Theorem~\ref{thm:main} with one potential and one critical point.
